\documentclass{article}
\usepackage{mathnotes}

\title{Rings and Fields}
\description{Algebraic structures with two operations - rings and fields including their axioms and fundamental properties.}

\begin{document}

\section{Rings}

\begin{definition}[Ring]
A \textbf{ring} is a set together with two binary operations $+$ and $\cdot$, which we will call addition and multiplication, such that the following axioms are satisfied:

\begin{enumerate}
\item $\langle R, + \rangle$ is an \@{Abelian} \@{group}.
\end{enumerate}

\begin{enumerate}
\item Multiplication is associative.
\end{enumerate}

\begin{enumerate}
\item For all $a,b,c \in R$, the left distributive law and the right distributive law hold, i.e.
\end{enumerate}

\[ a \cdot (b + c) = (a \cdot b) + (a \cdot c), \quad (a + b) \cdot c = (a \cdot c) + (b \cdot c). \]
\end{definition}

\begin{example}[Examples of Rings]
The integers, rationals, reals and complex numbers are all \@{rings} with the usual addition and multiplication.
\end{example}

\begin{definition}[Ring Homomorphism]
A \textbf{ring homomorphism} $\phi : R \to R'$ must satisfy the following two properties:

\begin{enumerate}
\item $\phi{(a+b)} = \phi{(a)} + \phi{(b)}.$ 
\end{enumerate}

\begin{enumerate}
\item $\phi{(ab)} = \phi{(a)}\phi{(b)}.$
\end{enumerate}
\end{definition}

\begin{remark}\label{ring-multiplicative-identity}
A ring doesn't have to have a \@[multiplicative identity element]{unity}, but it can. If it has one, it's denoted $1$ and for all $a$ in $R$ satisfies $a1 = 1a = a$.
\end{remark}

\begin{definition}[unity]
The element $1$ is also called \textbf{unity}.
\end{definition}

\begin{definition}[Commutative Ring]
A ring in which multiplication is commutative is called a \textbf{commutative ring}.
\end{definition}

\begin{definition}[Ring with Unity]
A ring that has a multiplicative identity element is called a \textbf{ring with unity.}
\end{definition}

\begin{definition}[Multiplicative Inverse]
For some element $a$ in a ring with unity $R$ where $1 \neq 0$, if $a^{-1} \in R$ such that $aa^{-1} = a^{-1}a = 1$, $a^{-1}$ is said to be the \textbf{multiplicative inverse} of $a.$
\end{definition}

\begin{definition}[Unit]
If $a$ has a multiplicative inverse in $R,$ $a$ is said to be a \textbf{unit} in $R$.
\end{definition}

\subsubsection{Fields}

\begin{definition}[Division Ring]
Let $R$ be a ring with unity. If every nonzero element of $R$ is a unit (has a multiplicative inverse), then $R$ is called a \textbf{division ring.}
\end{definition}

\begin{definition}[Field]
A commutative division ring is called a \textbf{field}.
\end{definition}

\begin{example}[Examples of Fields]
The integers are not a \@{field}, but the rationals, the reals, and the complex numbers are all fields.
\end{example}

\end{document}
